模型的目标函数、能量平衡、储电量递推和变量边界均为线性关系，因此将 144 时段变量和约束组装为稀疏矩阵，调用 HiGHS 求解。HiGHS 提供高性能线性规划算法\cite{highs}，求解接口可返回求解状态、原始残差和对偶信息。数值解中 $\max_t(c_td_t)=0$，说明本算例的连续 LP 解满足充放电互斥要求；其与显式互斥模型的对照结果在结果分析中给出。

\noindent\textbf{HiGHS 稀疏线性规划算法}
\algorithmstart
\algostep{1}{读取并校验数据}
按附件原始行序读取 144 个时段，核查电价、负荷和光伏数据的完整性与量纲。
\algostep{2}{完成单位换算}
将负荷功率和光伏功率乘以 $\Delta t=1/6\,\mathrm{h}$，得到各时段电量。
\algostep{3}{构造线性规划}
建立 $g,c,d,E,s$ 五类连续变量，写入目标函数、能量平衡、储电量递推和上下界。
\algostep{4}{调用 HiGHS 求解}
传入目标系数向量、稀疏约束矩阵与变量边界，得到最优购电计划和储能轨迹。
\algostep{5}{进行可行性审计}
检查平衡残差、储电量边界、首末储电量和同时充放电量；不通过时返回 Step 1 复核数据口径。
\algostep{6}{按题设表格回填}
把 144 个购电量逐行写入“计划购电量”，按连续 24 个时段汇总六个四小时区间，写入“充放电量”，并按题目表 1、表 2 的格式展示正文结果。
\algorithmend
